\documentclass[11pt,a4paper]{article}

% ---- Packages ----
\usepackage[utf8]{inputenc}
\usepackage[T1]{fontenc}
\usepackage{amsmath,amssymb,amsthm}
\usepackage{mathtools}
\usepackage[margin=1in,headheight=14pt]{geometry}
\usepackage{hyperref}
\usepackage{enumitem}
\usepackage{xcolor}
\usepackage{tcolorbox}
\usepackage{fancyhdr}
\usepackage{booktabs}

% ---- Hyperref ----
\hypersetup{
    colorlinks=true,
    linkcolor=red,
    citecolor=blue,
    urlcolor=blue
}

% ---- Header ----
\pagestyle{fancy}
\fancyhf{}
\fancyhead[L]{\textit{Lesson 11: Static and Extensive-Form Games}}
\fancyhead[R]{\thepage}
\renewcommand{\headrulewidth}{0.4pt}

% ---- Theorem Environments ----
\theoremstyle{definition}
\newtheorem{definition}{Definition}[section]
\newtheorem{example}{Example}[section]
\newtheorem{exercise}{Exercise}[section]

\theoremstyle{plain}
\newtheorem{theorem}{Theorem}[section]
\newtheorem{lemma}[theorem]{Lemma}
\newtheorem{proposition}[theorem]{Proposition}
\newtheorem{corollary}[theorem]{Corollary}

\theoremstyle{remark}
\newtheorem{remark}{Remark}[section]

% ---- Colored Boxes ----
\tcbuselibrary{theorems,skins,breakable}

\newtcolorbox{keyresult}[1][]{
    colback=green!5!white,
    colframe=green!50!black,
    fonttitle=\bfseries,
    title=#1,
    breakable
}

\newtcolorbox{rlconnection}[1][]{
    colback=blue!5!white,
    colframe=blue!50!black,
    fonttitle=\bfseries,
    title=#1,
    breakable
}

\newtcolorbox{warning}[1][]{
    colback=red!5!white,
    colframe=red!50!black,
    fonttitle=\bfseries,
    title=#1,
    breakable
}

% ---- Operators ----
\newcommand{\R}{\mathbb{R}}
\newcommand{\N}{\mathbb{N}}
\newcommand{\Q}{\mathbb{Q}}
\newcommand{\Z}{\mathbb{Z}}
\newcommand{\C}{\mathbb{C}}
\newcommand{\Prob}{\mathbb{P}}
\newcommand{\E}{\mathbb{E}}
\newcommand{\indicator}{\mathbf{1}}
\DeclareMathOperator{\supp}{supp}
\DeclareMathOperator{\conv}{conv}
\DeclareMathOperator{\argmax}{argmax}
\DeclareMathOperator{\argmin}{argmin}
\DeclareMathOperator{\BR}{BR}

% ---- Title ----
\title{\textbf{Lesson 11: Static and Extensive-Form Games}\\[10pt]
\large Nash Equilibrium, Minimax Theorem, and Subgame Perfection\\[5pt]
\normalsize Session 11 of 102 $\mid$ Phase I: Mathematical Foundations $\mid$ 8\,h\\
\normalsize Primary text: Fudenberg \& Tirole, \emph{Game Theory}, Chapters 1--3}
\author{Curriculum: A Rigorous Learning Path to Multi-Agent Deep RL}
\date{}

\begin{document}

\maketitle

\begin{abstract}
This lesson develops the foundational theory of strategic interaction: games in which
multiple decision-makers choose actions simultaneously or sequentially, each seeking to
maximize their own payoff. We define normal-form games, prove Nash's existence theorem
for mixed-strategy equilibria via Kakutani's fixed point theorem, establish
von Neumann's minimax theorem for zero-sum games, introduce correlated equilibrium,
and then pass to extensive-form games where we define subgame perfect equilibrium
and prove Zermelo's theorem by backward induction. Every theorem is proved in full.
The reader who masters this material will possess the game-theoretic language and
solution concepts that underpin all multi-agent reinforcement learning algorithms
in Phase~05.
\end{abstract}

\tableofcontents
\newpage

%% ============================================================
%% SECTION 1: INTRODUCTION
%% ============================================================
\section{Why Game Theory for RL and ML?}
\label{sec:intro}

In single-agent reinforcement learning, the environment is typically modeled as a
Markov decision process (MDP): the agent chooses actions, the environment responds
according to a fixed (if unknown) transition kernel, and the agent seeks a policy
maximizing cumulative reward. The key simplification is that the environment does
not \emph{adapt} to the agent's behavior.

Multi-agent reinforcement learning (MARL) breaks this assumption. When two or more
learning agents interact in a shared environment, each agent's ``environment'' includes
the other agents, whose behavior changes as they learn. The resulting non-stationarity
means that standard single-agent convergence guarantees fail. To analyze MARL, we need
a mathematical framework for strategic interaction --- and that framework is game theory.

\begin{rlconnection}[MARL Is Game Theory + Learning]
Every MARL algorithm implicitly assumes a game-theoretic solution concept:
\begin{itemize}[nosep]
    \item \textbf{Independent Q-learning}: each agent performs Q-learning ignoring
    others. If it converges, the result is a Nash equilibrium of the stage game
    (under restrictive conditions).
    \item \textbf{Minimax-Q} (Littman, 1994): assumes a two-player zero-sum
    stochastic game and seeks the minimax value --- directly applying
    von Neumann's theorem (Theorem~\ref{thm:minimax}).
    \item \textbf{Nash-Q} (Hu \& Wellman, 2003): computes Nash equilibria of
    stage games at each step.
    \item \textbf{MAPPO, QMIX}: cooperative MARL algorithms whose convergence
    analysis relies on the team game structure (common payoff games).
\end{itemize}
Without understanding Nash equilibria, minimax, and subgame perfection, one cannot
state what these algorithms are trying to compute, let alone prove convergence.
\end{rlconnection}

This lesson covers three progressively richer models of strategic interaction:
\begin{enumerate}[nosep]
    \item \textbf{Normal-form (static) games} (Sections~\ref{sec:normal_form}--\ref{sec:correlated}):
    players choose actions simultaneously.
    \item \textbf{Extensive-form (dynamic) games} (Sections~\ref{sec:extensive}--\ref{sec:spe}):
    players move sequentially, possibly with imperfect information.
    \item \textbf{Solution concepts}: Nash equilibrium, minimax strategies, correlated
    equilibrium, and subgame perfect equilibrium.
\end{enumerate}

\noindent\textbf{Prerequisites.} Familiarity with basic probability (finite probability
spaces suffice for this lesson), linear algebra (convex sets, hyperplane separation),
and real analysis (continuity, compactness in $\R^n$). Knowledge of linear programming
duality is helpful but not required; we state the needed result.

%% ============================================================
%% SECTION 2: NORMAL-FORM GAMES
%% ============================================================
\section{Normal-Form Games}
\label{sec:normal_form}

\subsection{Basic Definitions}
\label{subsec:normal_form_def}

\begin{definition}[Normal-Form Game]
\label{def:normal_form}
A \emph{normal-form game} (or \emph{strategic-form game}) is a tuple
$\Gamma = (N, (S_i)_{i \in N}, (u_i)_{i \in N})$ where:
\begin{enumerate}[nosep,label=(\roman*)]
    \item $N = \{1, 2, \ldots, n\}$ is a finite set of \emph{players}.
    \item For each $i \in N$, $S_i$ is a non-empty set of \emph{strategies}
    (or \emph{actions}) available to player~$i$. We write
    $S = \prod_{i \in N} S_i$ for the set of \emph{strategy profiles}.
    \item For each $i \in N$, $u_i : S \to \R$ is player~$i$'s \emph{payoff function}
    (or \emph{utility function}).
\end{enumerate}
The game is \emph{finite} if each $S_i$ is a finite set.
\end{definition}

\begin{remark}[Notation for ``All but $i$'']
\label{rem:notation}
For a strategy profile $s = (s_1, \ldots, s_n) \in S$, we write
$s_{-i} = (s_1, \ldots, s_{i-1}, s_{i+1}, \ldots, s_n) \in S_{-i}
:= \prod_{j \neq i} S_j$.
We often write $s = (s_i, s_{-i})$ to emphasize player~$i$'s strategy.
\end{remark}

\begin{definition}[Bimatrix Game]
\label{def:bimatrix}
A \emph{bimatrix game} is a two-player finite game presented as a pair of matrices
$(A, B) \in \R^{m \times n} \times \R^{m \times n}$, where player~1 has $m$ strategies,
player~2 has $n$ strategies, $A_{jk} = u_1(j,k)$, and $B_{jk} = u_2(j,k)$.
If $B = -A$, the game is \emph{zero-sum}.
\end{definition}

\subsection{Classical Examples}
\label{subsec:examples}

\begin{example}[Prisoner's Dilemma]
\label{ex:pd}
Two players simultaneously choose to \emph{Cooperate} (C) or \emph{Defect} (D):
\[
\begin{array}{c|cc}
 & C & D \\ \hline
C & (3,3) & (0,5) \\
D & (5,0) & (1,1)
\end{array}
\]
Here $D$ strictly dominates $C$ for both players (Definition~\ref{def:dominance}),
so the unique Nash equilibrium is $(D,D)$ with payoffs $(1,1)$, despite $(C,C)$
giving both players a higher payoff.
\end{example}

\begin{example}[Battle of the Sexes]
\label{ex:bos}
Players~1 and~2 prefer different activities ($O$pera vs.\ $F$ootball) but prefer
coordination to miscoordination:
\[
\begin{array}{c|cc}
 & O & F \\ \hline
O & (3,1) & (0,0) \\
F & (0,0) & (1,3)
\end{array}
\]
This game has two pure Nash equilibria --- $(O,O)$ and $(F,F)$ --- and one mixed
equilibrium (computed in Example~\ref{ex:mixed_bos}).
\end{example}

\begin{example}[Matching Pennies]
\label{ex:mp}
A zero-sum game: each player chooses $H$eads or $T$ails. Player~1 wins if
they match; player~2 wins if they differ:
\[
\begin{array}{c|cc}
 & H & T \\ \hline
H & (1,-1) & (-1,1) \\
T & (-1,1) & (1,-1)
\end{array}
\]
This game has no pure Nash equilibrium (Proposition~\ref{prop:mp_no_pure}),
motivating the extension to mixed strategies.
\end{example}

\begin{example}[Stag Hunt]
\label{ex:stag}
Players choose to hunt a $S$tag (requiring cooperation) or a $H$are (safe but
less rewarding):
\[
\begin{array}{c|cc}
 & S & H \\ \hline
S & (4,4) & (0,3) \\
H & (3,0) & (3,3)
\end{array}
\]
Both $(S,S)$ and $(H,H)$ are pure Nash equilibria. The equilibrium $(S,S)$ is
\emph{payoff-dominant} (higher payoffs for both), while $(H,H)$ is \emph{risk-dominant}
(safer against uncertainty about the opponent).
\end{example}

\begin{rlconnection}[Stag Hunt and Cooperative MARL]
The stag hunt captures a fundamental tension in cooperative MARL: agents can achieve
higher collective reward through coordination, but each agent bears risk if the others
fail to coordinate. This is precisely the challenge faced by CTDE (centralized training,
decentralized execution) algorithms like QMIX and MAPPO.
\end{rlconnection}

\subsection{Dominant and Dominated Strategies}
\label{subsec:dominance}

\begin{definition}[Dominance]
\label{def:dominance}
Let $\Gamma = (N, (S_i), (u_i))$ be a normal-form game and fix player~$i$.
A strategy $s_i \in S_i$ \emph{strictly dominates} $s_i' \in S_i$ if
\[
    u_i(s_i, s_{-i}) > u_i(s_i', s_{-i}) \quad \text{for all } s_{-i} \in S_{-i}.
\]
A strategy $s_i$ \emph{weakly dominates} $s_i'$ if
\[
    u_i(s_i, s_{-i}) \geq u_i(s_i', s_{-i}) \quad \text{for all } s_{-i} \in S_{-i},
\]
with strict inequality for at least one $s_{-i}$.
A strategy is \emph{strictly dominated} if some other strategy strictly dominates it.
A strategy is a \emph{(strictly) dominant strategy} if it strictly dominates every
other strategy of that player.
\end{definition}

\begin{proposition}[Dominated Strategies Are Never Best Responses]
\label{prop:dominated_not_br}
If $s_i'$ is strictly dominated by $s_i$, then $s_i'$ is never a best response:
for every $s_{-i} \in S_{-i}$, $s_i' \notin \argmax_{t_i \in S_i} u_i(t_i, s_{-i})$.
\end{proposition}

\begin{proof}
Fix any $s_{-i} \in S_{-i}$. Since $s_i$ strictly dominates $s_i'$, we have
$u_i(s_i, s_{-i}) > u_i(s_i', s_{-i})$, so $s_i'$ does not maximize $u_i(\cdot, s_{-i})$.
\end{proof}

\subsection{Iterated Elimination of Strictly Dominated Strategies}
\label{subsec:iesds}

\begin{definition}[IESDS]
\label{def:iesds}
\emph{Iterated elimination of strictly dominated strategies} (IESDS) is the following
procedure. Set $S_i^0 = S_i$ for all~$i$. For $k = 0, 1, 2, \ldots$, define
\[
    S_i^{k+1} = \big\{s_i \in S_i^k : \text{there is no } s_i' \in S_i^k
    \text{ such that } u_i(s_i', s_{-i}) > u_i(s_i, s_{-i})
    \text{ for all } s_{-i} \in S_{-i}^k\big\}.
\]
The set of \emph{surviving strategies} for player~$i$ is
$S_i^\infty = \bigcap_{k=0}^\infty S_i^k$.
\end{definition}

\begin{theorem}[Order Independence of IESDS]
\label{thm:iesds_order}
In a finite game, the set of strategy profiles surviving iterated elimination of
strictly dominated strategies is independent of the order of elimination.
\end{theorem}

\begin{proof}
We prove that any strategy eliminated in some order of elimination is eventually
eliminated in every order.

Let $s_i^*$ be a strategy of player~$i$ that is eliminated at round~$k$ in some
particular elimination order $E$. This means there exists $s_i' \in S_i^{k-1}(E)$
such that $u_i(s_i', s_{-i}) > u_i(s_i^*, s_{-i})$ for all $s_{-i} \in S_{-i}^{k-1}(E)$.

Consider an alternative elimination order $E'$. We show $s_i^*$ is eventually
eliminated under $E'$. Let $T = S_{-i}^\infty(E')$ be the surviving opponents'
strategies under $E'$.

Since $S_{-i}^\infty(E') \subseteq S_{-i}^{k-1}(E')$ and every strategy surviving
under any order is in particular not eliminated in the first $k-1$ rounds under $E$
(we prove this by induction on rounds), we need a cleaner argument.

We proceed by induction on the round of elimination. For the base case ($k=0$),
$s_i^*$ is strictly dominated in the original game, so it will be eliminated in
round~0 under any order.

For the inductive step, suppose every strategy eliminated in rounds $0, \ldots, k-1$
under \emph{any} elimination order is eventually eliminated under every order. Now
suppose $s_i^*$ is eliminated at round~$k$ under order~$E$. Then there exists
$s_i' \in S_i^{k-1}(E)$ with $u_i(s_i', s_{-i}) > u_i(s_i^*, s_{-i})$ for all
$s_{-i} \in S_{-i}^{k-1}(E)$.

Under order $E'$, the set $S_{-i}^{k-1}(E)$ may differ from $S_{-i}^{k-1}(E')$.
However, the key observation is:

\emph{Claim}: $S_{-i}^\infty(E') \subseteq S_{-i}^{k-1}(E)$.

To prove the claim, note that any strategy $s_j$ removed before round~$k$ under
order~$E$ was strictly dominated by some $s_j'$ against strategies surviving at
that point. By the inductive hypothesis, $s_j$ is eventually removed under $E'$ as
well, so $s_j \notin S_j^\infty(E')$. Taking contrapositives:
$s_j \in S_j^\infty(E') \implies s_j \in S_j^{k-1}(E)$, establishing the claim.

Now, since $S_{-i}^\infty(E') \subseteq S_{-i}^{k-1}(E)$, the domination
$u_i(s_i', s_{-i}) > u_i(s_i^*, s_{-i})$ holds for all $s_{-i} \in S_{-i}^\infty(E')$.
If $s_i' \in S_i^\infty(E')$, then $s_i^*$ is dominated (by $s_i'$) relative to
the surviving strategies under $E'$ and hence is eliminated. If $s_i'$ is itself
eliminated under $E'$, then $s_i'$ was dominated by some $s_i''$, and by transitivity
of strict domination on the surviving set, $s_i^*$ remains dominated and is eliminated.

By induction, the surviving sets coincide for all elimination orders.
\end{proof}

\subsection{Rationalizability}
\label{subsec:rationalizability}

\begin{definition}[Best Response]
\label{def:best_response}
Player~$i$'s \emph{best response correspondence} is the set-valued map
$\BR_i : S_{-i} \rightrightarrows S_i$ defined by
\[
    \BR_i(s_{-i}) = \argmax_{s_i \in S_i} u_i(s_i, s_{-i}).
\]
\end{definition}

\begin{definition}[Rationalizable Strategies]
\label{def:rationalizable}
A strategy $s_i$ is \emph{rationalizable} if it survives the iterated elimination
of strategies that are never best responses. Formally, set $R_i^0 = S_i$ and define
\[
    R_i^{k+1} = \big\{s_i \in R_i^k : \exists\, s_{-i} \in R_{-i}^k
    \text{ such that } s_i \in \BR_i(s_{-i})\big\}.
\]
The set of rationalizable strategies is $R_i^\infty = \bigcap_{k=0}^\infty R_i^k$.
\end{definition}

\begin{proposition}[IESDS and Rationalizability in Two-Player Games]
\label{prop:iesds_rational}
In finite two-player games, the set of rationalizable strategies coincides with the
set of strategies surviving IESDS.
\end{proposition}

\begin{proof}
Every strictly dominated strategy is never a best response
(Proposition~\ref{prop:dominated_not_br}), so
$S_i^\infty \supseteq R_i^\infty$. Conversely, in a two-player finite game,
a strategy that is never a best response to any pure strategy of the opponent is
strictly dominated (this follows from the supporting hyperplane theorem applied
to the payoff polyhedron). Hence a strategy eliminated in the rationalizability
procedure is strictly dominated relative to the surviving set, giving
$R_i^\infty \supseteq S_i^\infty$.
\end{proof}

\begin{warning}[Three or More Players]
For games with three or more players, IESDS may not eliminate all never-best-response
strategies. Rationalizability can be a strictly coarser concept than IESDS survival
in games with $n \geq 3$ players when we restrict to pure strategies. The equivalence
is restored when we extend dominance to allow domination by mixed strategies.
\end{warning}

%% ============================================================
%% SECTION 3: NASH EQUILIBRIUM
%% ============================================================
\section{Nash Equilibrium in Pure Strategies}
\label{sec:nash_pure}

\begin{definition}[Pure Strategy Nash Equilibrium]
\label{def:nash_pure}
A strategy profile $s^* = (s_1^*, \ldots, s_n^*) \in S$ is a \emph{(pure strategy)
Nash equilibrium} (NE) if no player can unilaterally improve their payoff:
\[
    u_i(s_i^*, s_{-i}^*) \geq u_i(s_i, s_{-i}^*)
    \quad \text{for all } s_i \in S_i, \text{ for all } i \in N.
\]
Equivalently, $s_i^* \in \BR_i(s_{-i}^*)$ for all $i \in N$.
\end{definition}

\begin{proposition}[NE as Mutual Best Response]
\label{prop:ne_br}
A strategy profile $s^*$ is a Nash equilibrium if and only if
$s_i^* \in \BR_i(s_{-i}^*)$ for every $i \in N$.
\end{proposition}

\begin{proof}
By Definition~\ref{def:nash_pure}, $s^*$ is a NE if and only if
$u_i(s_i^*, s_{-i}^*) \geq u_i(s_i, s_{-i}^*)$ for all $s_i \in S_i$ and all~$i$.
By Definition~\ref{def:best_response}, this is precisely
$s_i^* \in \argmax_{s_i \in S_i} u_i(s_i, s_{-i}^*) = \BR_i(s_{-i}^*)$ for all~$i$.
\end{proof}

\begin{proposition}[Matching Pennies Has No Pure NE]
\label{prop:mp_no_pure}
The Matching Pennies game (Example~\ref{ex:mp}) has no pure strategy Nash equilibrium.
\end{proposition}

\begin{proof}
We check all four profiles:
\begin{itemize}[nosep]
    \item $(H,H)$: Player~2 gains by deviating to~$T$ (payoff $-1 \to 1$).
    \item $(H,T)$: Player~1 gains by deviating to~$T$ (payoff $-1 \to 1$).
    \item $(T,H)$: Player~1 gains by deviating to~$H$ (payoff $-1 \to 1$).
    \item $(T,T)$: Player~2 gains by deviating to~$H$ (payoff $-1 \to 1$).
\end{itemize}
Every profile admits a profitable deviation, so no pure NE exists.
\end{proof}

\begin{remark}
The Prisoner's Dilemma has a unique NE at $(D,D)$, the Battle of the Sexes has two
pure NE at $(O,O)$ and $(F,F)$, and the Stag Hunt has two at $(S,S)$ and $(H,H)$.
Matching Pennies has none --- motivating the extension to mixed strategies.
\end{remark}

%% ============================================================
%% SECTION 4: MIXED STRATEGIES
%% ============================================================
\section{Mixed Strategies and Mixed Extensions}
\label{sec:mixed}

\subsection{Definitions}
\label{subsec:mixed_def}

\begin{definition}[Mixed Strategy]
\label{def:mixed_strategy}
Let $\Gamma$ be a finite game with strategy set $S_i = \{s_i^1, \ldots, s_i^{m_i}\}$
for player~$i$. A \emph{mixed strategy} for player~$i$ is a probability distribution
$\sigma_i \in \Delta(S_i)$, where
\[
    \Delta(S_i) = \bigg\{\sigma_i \in \R^{m_i} : \sigma_i(s_i^j) \geq 0
    \text{ for all } j, \;\; \sum_{j=1}^{m_i} \sigma_i(s_i^j) = 1\bigg\}
\]
is the \emph{$(m_i-1)$-simplex} over $S_i$. A \emph{mixed strategy profile} is
$\sigma = (\sigma_1, \ldots, \sigma_n) \in \Delta := \prod_{i \in N} \Delta(S_i)$.
\end{definition}

\begin{definition}[Support]
\label{def:support}
The \emph{support} of a mixed strategy $\sigma_i$ is
$\supp(\sigma_i) = \{s_i \in S_i : \sigma_i(s_i) > 0\}$.
A pure strategy $s_i$ is identified with the degenerate distribution placing
probability~1 on $s_i$.
\end{definition}

\begin{definition}[Mixed Extension]
\label{def:mixed_extension}
The \emph{mixed extension} of a finite game $\Gamma = (N, (S_i), (u_i))$ is the game
$\bar{\Gamma} = (N, (\Delta(S_i)), (\bar{u}_i))$ where the expected payoff under
independent mixing is
\[
    \bar{u}_i(\sigma) = \sum_{s \in S} \bigg(\prod_{j \in N} \sigma_j(s_j)\bigg) u_i(s)
    = \E_{s \sim \sigma}[u_i(s)].
\]
We drop the bar and write $u_i(\sigma)$ when the context is clear.
\end{definition}

\begin{remark}[Multilinearity]
\label{rem:multilinear}
The expected payoff $u_i(\sigma)$ is \emph{multilinear}: it is linear in each
$\sigma_i$ when the other players' strategies are held fixed. Explicitly,
\[
    u_i(\sigma_i, \sigma_{-i}) = \sum_{s_i \in S_i} \sigma_i(s_i)
    \underbrace{\sum_{s_{-i} \in S_{-i}}
    \bigg(\prod_{j \neq i} \sigma_j(s_j)\bigg) u_i(s_i, s_{-i})}_{=:\, u_i(s_i, \sigma_{-i})}.
\]
The quantity $u_i(s_i, \sigma_{-i})$ is the expected payoff to player~$i$ from playing
the pure strategy $s_i$ when opponents play $\sigma_{-i}$.
\end{remark}

\subsection{Mixed Nash Equilibrium and the Indifference Principle}
\label{subsec:mixed_ne}

\begin{definition}[Mixed Strategy Nash Equilibrium]
\label{def:mixed_ne}
A mixed strategy profile $\sigma^*$ is a \emph{Nash equilibrium} if
\[
    u_i(\sigma_i^*, \sigma_{-i}^*) \geq u_i(\sigma_i, \sigma_{-i}^*)
    \quad \text{for all } \sigma_i \in \Delta(S_i), \text{ for all } i \in N.
\]
\end{definition}

\begin{theorem}[Indifference Principle]
\label{thm:indifference}
Let $\sigma^*$ be a mixed strategy Nash equilibrium. For each player~$i$, every pure
strategy in $\supp(\sigma_i^*)$ yields the same expected payoff against $\sigma_{-i}^*$,
and this common payoff is at least as large as the expected payoff from any pure
strategy outside the support. Formally:
\begin{enumerate}[nosep,label=(\alph*)]
    \item If $s_i, s_i' \in \supp(\sigma_i^*)$, then
    $u_i(s_i, \sigma_{-i}^*) = u_i(s_i', \sigma_{-i}^*)$.
    \item If $s_i \in \supp(\sigma_i^*)$ and $t_i \notin \supp(\sigma_i^*)$, then
    $u_i(s_i, \sigma_{-i}^*) \geq u_i(t_i, \sigma_{-i}^*)$.
\end{enumerate}
\end{theorem}

\begin{proof}
By multilinearity (Remark~\ref{rem:multilinear}),
\[
    u_i(\sigma_i^*, \sigma_{-i}^*) = \sum_{s_i \in S_i} \sigma_i^*(s_i)\, u_i(s_i, \sigma_{-i}^*).
\]
Let $v_i^* = u_i(\sigma_i^*, \sigma_{-i}^*)$ be the equilibrium payoff for player~$i$.

\textbf{Part (b):} Suppose $t_i \in S_i$ satisfies $u_i(t_i, \sigma_{-i}^*) > v_i^*$.
Then the pure strategy $t_i$ (viewed as a degenerate mixed strategy) gives payoff
$u_i(t_i, \sigma_{-i}^*) > v_i^* = u_i(\sigma_i^*, \sigma_{-i}^*)$, contradicting the
NE condition. Hence $u_i(s_i, \sigma_{-i}^*) \leq v_i^*$ for all $s_i$ with equality
impossible for all $s_i$ unless we also have a lower bound. Since
$v_i^*$ is a convex combination of the payoffs $u_i(s_i, \sigma_{-i}^*)$ with
weights $\sigma_i^*(s_i) \geq 0$ summing to~1, if any payoff in the support were
strictly less than $v_i^*$, then some payoff (in or outside the support) must be
strictly greater than $v_i^*$ (to make the weighted average equal $v_i^*$), which
we just showed is impossible. Therefore:
\begin{itemize}[nosep]
    \item For all $s_i \in \supp(\sigma_i^*)$: $u_i(s_i, \sigma_{-i}^*) = v_i^*$.
    \item For all $s_i \notin \supp(\sigma_i^*)$: $u_i(s_i, \sigma_{-i}^*) \leq v_i^*$.
\end{itemize}

\textbf{Part (a):} Immediate from the first bullet point.
\end{proof}

\begin{example}[Mixed NE in Battle of the Sexes]
\label{ex:mixed_bos}
In the Battle of the Sexes (Example~\ref{ex:bos}), let player~1 play $O$ with
probability $p$ and player~2 play $O$ with probability $q$. By the indifference
principle, in a fully mixed NE, player~1 must be indifferent:
\[
    u_1(O, q) = u_1(F, q) \implies 3q + 0(1-q) = 0 \cdot q + 1 \cdot (1-q)
    \implies 3q = 1 - q \implies q = \tfrac{1}{4}.
\]
Similarly, player~2 must be indifferent:
\[
    u_2(p, O) = u_2(p, F) \implies 1 \cdot p + 0(1-p) = 0 \cdot p + 3(1-p)
    \implies p = 3 - 3p \implies p = \tfrac{3}{4}.
\]
The mixed NE is $\sigma_1 = (\tfrac{3}{4}, \tfrac{1}{4})$,
$\sigma_2 = (\tfrac{1}{4}, \tfrac{3}{4})$, with expected payoffs
$(3 \cdot \tfrac{1}{4}, 1 \cdot \tfrac{3}{4}) = (\tfrac{3}{4}, \tfrac{3}{4})$.
Both players fare worse in the mixed NE than in either pure NE.
\end{example}

\begin{rlconnection}[Mixed Strategies and Stochastic Policies]
In RL, a stochastic policy $\pi(a|s)$ assigns probabilities to actions given a state.
This is precisely a mixed strategy in the stage game at state~$s$. The indifference
principle has a direct analog: in an optimal stochastic policy, every action in the
support must yield the same Q-value. If some action had a strictly higher Q-value,
the agent should place all probability on it. This is why $\epsilon$-greedy exploration
is suboptimal --- it assigns probability to actions that may not be in the support of
any Nash equilibrium.
\end{rlconnection}

%% ============================================================
%% SECTION 5: NASH EXISTENCE THEOREM
%% ============================================================
\section{Nash's Existence Theorem}
\label{sec:nash_existence}

The central result of this lesson: every finite game has at least one Nash equilibrium
in mixed strategies. The proof uses Kakutani's fixed point theorem.

\subsection{Kakutani's Fixed Point Theorem}
\label{subsec:kakutani}

\begin{definition}[Upper Hemicontinuity]
\label{def:uhc}
Let $X \subseteq \R^m$ and $Y \subseteq \R^n$. A correspondence (set-valued map)
$F : X \rightrightarrows Y$ is \emph{upper hemicontinuous} (uhc) at $x_0 \in X$ if
for every open set $V \supseteq F(x_0)$, there exists an open set $U \ni x_0$ such
that $F(x) \subseteq V$ for all $x \in U \cap X$.

Equivalently (when $Y$ is compact), $F$ is uhc at $x_0$ if whenever $x_k \to x_0$
and $y_k \in F(x_k)$ with $y_k \to y_0$, we have $y_0 \in F(x_0)$. This is the
\emph{closed graph} characterization.
\end{definition}

\begin{theorem}[Kakutani's Fixed Point Theorem]
\label{thm:kakutani}
Let $K \subseteq \R^n$ be non-empty, compact, and convex. Let $F : K \rightrightarrows K$
be a correspondence such that:
\begin{enumerate}[nosep,label=(\roman*)]
    \item $F(x)$ is non-empty for every $x \in K$.
    \item $F(x)$ is convex for every $x \in K$.
    \item $F$ is upper hemicontinuous.
\end{enumerate}
Then $F$ has a fixed point: there exists $x^* \in K$ with $x^* \in F(x^*)$.
\end{theorem}

\begin{proof}
We provide the standard proof via Brouwer's fixed point theorem.

For each $k \in \N$, let $T_k$ be a triangulation of $K$ with mesh size at most
$1/k$ (the maximum diameter of any simplex in the triangulation). For each vertex $v$
of $T_k$, choose $f_k(v) \in F(v)$ (possible since $F(v)$ is non-empty). Extend
$f_k$ to all of $K$ by convex interpolation on each simplex: if
$x = \sum_{j} \lambda_j v_j$ is the barycentric representation of $x$ in its simplex,
set $f_k(x) = \sum_j \lambda_j f_k(v_j)$.

We claim each $f_k : K \to K$ is continuous. Continuity within each simplex follows
from the linearity of the interpolation. On shared faces of adjacent simplices, the
interpolation agrees because barycentric coordinates are unique on the shared face.
Since $K$ is convex and each $f_k(x)$ is a convex combination of points in $K$,
we have $f_k(x) \in K$.

By Brouwer's fixed point theorem (applicable since $K$ is compact and convex, and
$f_k$ is continuous), each $f_k$ has a fixed point $x_k \in K$ with $f_k(x_k) = x_k$.

Since $K$ is compact, the sequence $(x_k)$ has a convergent subsequence
$x_{k_l} \to x^* \in K$.

We show $x^* \in F(x^*)$. For each $l$, let $\sigma_l$ be the simplex of $T_{k_l}$
containing $x_{k_l}$, with vertices $v_1^l, \ldots, v_p^l$ (where $p \leq n+1$).
Then
\[
    x_{k_l} = f_{k_l}(x_{k_l}) = \sum_{j=1}^p \lambda_j^l f_{k_l}(v_j^l),
    \quad \text{where } x_{k_l} = \sum_{j=1}^p \lambda_j^l v_j^l.
\]
Since the mesh size tends to~0, $\|v_j^l - x_{k_l}\| \leq 1/k_l \to 0$, so
$v_j^l \to x^*$ for each~$j$.

Now, $f_{k_l}(v_j^l) \in F(v_j^l)$. Since $K$ is compact, by passing to a further
subsequence (using a diagonal argument over the finitely many vertices), we may assume
$f_{k_l}(v_j^l) \to y_j$ for some $y_j \in K$.

By the closed-graph characterization of upper hemicontinuity (Definition~\ref{def:uhc}):
$v_j^l \to x^*$ and $f_{k_l}(v_j^l) \to y_j$ with $f_{k_l}(v_j^l) \in F(v_j^l)$
implies $y_j \in F(x^*)$.

The barycentric coordinates $\lambda_j^l \in [0,1]$ lie in a compact set, so
(passing to a further subsequence) $\lambda_j^l \to \lambda_j^*$ with
$\sum_j \lambda_j^* = 1$ and $\lambda_j^* \geq 0$.

Taking the limit in $x_{k_l} = \sum_j \lambda_j^l f_{k_l}(v_j^l)$:
\[
    x^* = \sum_j \lambda_j^* y_j.
\]
Since each $y_j \in F(x^*)$ and $F(x^*)$ is convex, the convex combination
$x^* = \sum_j \lambda_j^* y_j \in F(x^*)$.
\end{proof}

\subsection{Nash's Theorem}
\label{subsec:nash_theorem}

\begin{keyresult}[Nash's Existence Theorem]
\begin{theorem}[Nash, 1950]
\label{thm:nash_existence}
Every finite normal-form game has at least one mixed strategy Nash equilibrium.
\end{theorem}
\end{keyresult}

\begin{proof}
Let $\Gamma = (N, (S_i)_{i \in N}, (u_i)_{i \in N})$ be a finite game with
$|S_i| = m_i$. Define the set of mixed strategy profiles
$\Delta = \prod_{i \in N} \Delta(S_i)$, which is a product of simplices and hence
compact and convex in $\R^{m_1 + \cdots + m_n}$.

\textbf{Step 1: Define the best response correspondence.}
For each player~$i$, the \emph{mixed best response correspondence}
$\BR_i : \Delta_{-i} \rightrightarrows \Delta(S_i)$ is
\[
    \BR_i(\sigma_{-i}) = \argmax_{\sigma_i \in \Delta(S_i)} u_i(\sigma_i, \sigma_{-i}).
\]
Define the joint best response correspondence $\BR : \Delta \rightrightarrows \Delta$ by
\[
    \BR(\sigma) = \BR_1(\sigma_{-1}) \times \BR_2(\sigma_{-2}) \times \cdots
    \times \BR_n(\sigma_{-n}).
\]

\textbf{Step 2: $\BR(\sigma)$ is non-empty.}
For each~$i$, the function $\sigma_i \mapsto u_i(\sigma_i, \sigma_{-i})$ is
continuous (in fact, linear) on the compact set $\Delta(S_i)$. By the extreme value
theorem, the maximum is attained, so $\BR_i(\sigma_{-i}) \neq \emptyset$.

\textbf{Step 3: $\BR(\sigma)$ is convex.}
Fix~$i$ and $\sigma_{-i}$. Let $\sigma_i, \sigma_i' \in \BR_i(\sigma_{-i})$ and
$\lambda \in [0,1]$. Set $\hat{\sigma}_i = \lambda \sigma_i + (1-\lambda)\sigma_i'$.
By linearity of $u_i$ in $\sigma_i$ (Remark~\ref{rem:multilinear}):
\[
    u_i(\hat{\sigma}_i, \sigma_{-i})
    = \lambda\, u_i(\sigma_i, \sigma_{-i}) + (1-\lambda)\, u_i(\sigma_i', \sigma_{-i})
    = \lambda\, v_i^* + (1-\lambda)\, v_i^* = v_i^*,
\]
where $v_i^* = \max_{\tau_i} u_i(\tau_i, \sigma_{-i})$. So $\hat{\sigma}_i$ also
achieves the maximum, hence $\hat{\sigma}_i \in \BR_i(\sigma_{-i})$. Thus
$\BR_i(\sigma_{-i})$ is convex. Since $\BR(\sigma)$ is a product of convex sets,
it is convex.

\textbf{Step 4: $\BR$ is upper hemicontinuous.}
We use the closed-graph characterization. Let $\sigma^k \to \sigma^*$ in $\Delta$
and $\tau^k \in \BR(\sigma^k)$ with $\tau^k \to \tau^*$. We must show
$\tau^* \in \BR(\sigma^*)$, i.e., $\tau_i^* \in \BR_i(\sigma_{-i}^*)$ for each~$i$.

Fix player~$i$. Since $\tau_i^k \in \BR_i(\sigma_{-i}^k)$, for every
$\sigma_i \in \Delta(S_i)$:
\[
    u_i(\tau_i^k, \sigma_{-i}^k) \geq u_i(\sigma_i, \sigma_{-i}^k).
\]
Since $u_i$ is continuous (it is multilinear on the compact set $\Delta$), taking
$k \to \infty$:
\[
    u_i(\tau_i^*, \sigma_{-i}^*) \geq u_i(\sigma_i, \sigma_{-i}^*).
\]
Since this holds for all $\sigma_i$, we have $\tau_i^* \in \BR_i(\sigma_{-i}^*)$.

\textbf{Step 5: Apply Kakutani.}
The correspondence $\BR : \Delta \rightrightarrows \Delta$ satisfies all hypotheses of
Kakutani's theorem (Theorem~\ref{thm:kakutani}): $\Delta$ is non-empty, compact, and
convex; $\BR(\sigma)$ is non-empty and convex for each $\sigma$; and $\BR$ is upper
hemicontinuous. Therefore, there exists $\sigma^* \in \Delta$ with
$\sigma^* \in \BR(\sigma^*)$.

By definition of $\BR$, $\sigma_i^* \in \BR_i(\sigma_{-i}^*)$ for all~$i$,
which by Proposition~\ref{prop:ne_br} (extended to mixed strategies) means $\sigma^*$
is a Nash equilibrium.
\end{proof}

\begin{remark}[Alternative Proof via Brouwer]
\label{rem:brouwer_proof}
An alternative proof avoids Kakutani by constructing a continuous function whose
fixed points are NE. For each player~$i$ and pure strategy $s_i^j$, define the
``excess payoff'':
\[
    g_i^j(\sigma) = \max\big\{0,\; u_i(s_i^j, \sigma_{-i}) - u_i(\sigma_i, \sigma_{-i})\big\}.
\]
Then define $f : \Delta \to \Delta$ by
\[
    f_i^j(\sigma) = \frac{\sigma_i(s_i^j) + g_i^j(\sigma)}
    {1 + \sum_{k=1}^{m_i} g_i^k(\sigma)}.
\]
One verifies: (i) $f$ maps $\Delta$ to $\Delta$; (ii) $f$ is continuous; (iii)
$\sigma^*$ is a fixed point of $f$ if and only if $g_i^j(\sigma^*) = 0$ for all
$i, j$, which holds if and only if $\sigma^*$ is a NE. Brouwer's theorem gives a
fixed point, hence a NE.
\end{remark}

\begin{rlconnection}[Existence Guarantees in MARL]
Nash's theorem guarantees that every finite MARL problem (modeled as a stochastic game
with finite state and action spaces) has at least one Nash equilibrium --- at least in
the stage games. This justifies the goal of Nash-Q learning (Hu \& Wellman, 2003),
which attempts to compute Nash equilibria at each state. However, the theorem is
non-constructive: it guarantees existence but says nothing about efficient computation.
Indeed, computing a Nash equilibrium is PPAD-complete (Daskalakis, Goldberg, Papadimitriou,
2009), meaning no polynomial-time algorithm is known.
\end{rlconnection}

%% ============================================================
%% SECTION 6: MINIMAX THEOREM
%% ============================================================
\section{Von Neumann's Minimax Theorem}
\label{sec:minimax}

\subsection{Two-Player Zero-Sum Games}
\label{subsec:zero_sum}

\begin{definition}[Two-Player Zero-Sum Game]
\label{def:zero_sum}
A two-player game $(N, (S_1, S_2), (u_1, u_2))$ is \emph{zero-sum} if
$u_1(s) + u_2(s) = 0$ for all $s \in S$. It is specified by a single payoff matrix
$A \in \R^{m \times n}$ where $u_1(i,j) = A_{ij}$ and $u_2(i,j) = -A_{ij}$.
Player~1 (the ``row player'') maximizes; player~2 (the ``column player'') minimizes.
\end{definition}

\begin{definition}[Maximin and Minimax Values]
\label{def:maximin}
For a zero-sum game with payoff matrix $A$, define:
\begin{align*}
    \underline{v} &= \max_{x \in \Delta_m} \min_{y \in \Delta_n} x^\top A y
    \quad \text{(maximin value: player~1's guarantee)}, \\
    \overline{v} &= \min_{y \in \Delta_n} \max_{x \in \Delta_m} x^\top A y
    \quad \text{(minimax value: player~2's guarantee)},
\end{align*}
where $\Delta_k$ denotes the $(k-1)$-simplex in $\R^k$.
\end{definition}

\begin{lemma}[Maximin $\leq$ Minimax]
\label{lem:maximin_leq}
For any matrix $A \in \R^{m \times n}$, $\underline{v} \leq \overline{v}$.
\end{lemma}

\begin{proof}
For any fixed $x \in \Delta_m$ and $y \in \Delta_n$:
\[
    \min_{y' \in \Delta_n} x^\top A y' \leq x^\top A y
    \leq \max_{x' \in \Delta_m} (x')^\top A y.
\]
Taking $\max$ over $x$ on the left and $\min$ over $y$ on the right:
\[
    \underline{v} = \max_x \min_{y'} x^\top A y'
    \leq \min_y \max_{x'} (x')^\top A y = \overline{v}. \qedhere
\]
\end{proof}

\begin{keyresult}[Von Neumann's Minimax Theorem]
\begin{theorem}[Von Neumann, 1928]
\label{thm:minimax}
For every matrix $A \in \R^{m \times n}$,
\[
    \max_{x \in \Delta_m} \min_{y \in \Delta_n} x^\top A y
    = \min_{y \in \Delta_n} \max_{x \in \Delta_m} x^\top A y.
\]
The common value $v^* = \underline{v} = \overline{v}$ is called the \emph{value}
of the game.
\end{theorem}
\end{keyresult}

\begin{proof}
We give a proof using LP duality. By Lemma~\ref{lem:maximin_leq}, it suffices to
show $\underline{v} \geq \overline{v}$.

\textbf{Step 1: Reformulate as linear programs.}
Player~1's maximin problem is:
\[
    \underline{v} = \max_{x \in \Delta_m} \min_{y \in \Delta_n} x^\top A y.
\]
Since $y \mapsto x^\top A y$ is linear in $y$ and $\Delta_n$ is a simplex, the minimum
over $y$ is attained at a vertex $e_j$ (the $j$-th standard basis vector), giving
$\min_y x^\top A y = \min_{j=1,\ldots,n} (A^\top x)_j$.

Thus $\underline{v} = \max_{x \in \Delta_m} \min_{j} (A^\top x)_j$. Introducing an
auxiliary variable $v$, this is equivalent to the linear program:
\begin{align}
    \text{(P)} \quad \max\;&\; v \notag \\
    \text{s.t.}\;&\; A^\top x \geq v\, \mathbf{1}, \quad x \in \Delta_m, \quad v \in \R,
    \label{eq:primal}
\end{align}
where $\mathbf{1} = (1, \ldots, 1)^\top \in \R^n$ and the inequality is componentwise.

\textbf{Step 2: The dual LP.}
Similarly, player~2's minimax problem gives:
\[
    \overline{v} = \min_{y \in \Delta_n} \max_{i} (Ay)_i.
\]
This is equivalent to:
\begin{align}
    \text{(D)} \quad \min\;&\; w \notag \\
    \text{s.t.}\;&\; Ay \leq w\, \mathbf{1}, \quad y \in \Delta_n, \quad w \in \R.
    \label{eq:dual}
\end{align}

\textbf{Step 3: LP duality.}
We show that (P) and (D) are LP duals of each other. Rewrite (P) in standard form.
Let $v = \underline{v}$. The constraints of (P) are:
$A^\top x - v\, \mathbf{1} \geq 0$, $\mathbf{1}^\top x = 1$, $x \geq 0$.
This is a linear program in the variables $(x, v) \in \R^{m+1}$.

The dual of this LP (after careful bookkeeping of the dual variables associated
with each primal constraint) yields precisely the LP (D). By strong LP duality
(both problems are feasible --- take $x = (1/m, \ldots, 1/m)^\top$ and
$y = (1/n, \ldots, 1/n)^\top$ --- and bounded), the optimal values are equal:
\[
    \underline{v} = \text{opt}(\text{P}) = \text{opt}(\text{D}) = \overline{v}. \qedhere
\]
\end{proof}

\begin{corollary}[NE in Zero-Sum Games]
\label{cor:ne_zero_sum}
In a two-player zero-sum game with payoff matrix $A$ and value $v^*$:
\begin{enumerate}[nosep,label=(\alph*)]
    \item $(x^*, y^*)$ is a Nash equilibrium if and only if $x^*$ is a maximin
    strategy and $y^*$ is a minimax strategy.
    \item All NE yield the same payoff: $u_1 = v^*$, $u_2 = -v^*$.
    \item NE strategies are interchangeable: if $(x^*, y^*)$ and $(x^{**}, y^{**})$
    are both NE, then so are $(x^*, y^{**})$ and $(x^{**}, y^*)$.
\end{enumerate}
\end{corollary}

\begin{proof}
\textbf{(a)} $(\Rightarrow)$: Let $(x^*, y^*)$ be a NE. Then $x^*$ is a best response
to $y^*$, so $(x^*)^\top A y^* = \max_x x^\top A y^*$. Similarly, $y^*$ minimizes
$x^\top A y$ against $x^*$, so $(x^*)^\top A y^* = \min_y (x^*)^\top A y$. Therefore:
\[
    \max_x x^\top A y^* = (x^*)^\top A y^* = \min_y (x^*)^\top A y.
\]
But $\min_y (x^*)^\top A y \leq \underline{v} \leq \overline{v}
\leq \max_x x^\top A y^*$,
and by the minimax theorem $\underline{v} = \overline{v} = v^*$, so all inequalities
are equalities. Hence $\min_y (x^*)^\top A y = v^*$ (so $x^*$ achieves the maximin)
and $\max_x x^\top A y^* = v^*$ (so $y^*$ achieves the minimax).

$(\Leftarrow)$: If $x^*$ achieves the maximin and $y^*$ achieves the minimax, then
\[
    (x^*)^\top A y^* \geq \min_y (x^*)^\top A y = v^*
    = \max_x x^\top A y^* \geq (x^*)^\top A y^*,
\]
so $(x^*)^\top A y^* = v^*$. For any $x$: $x^\top A y^* \leq \max_{x'} (x')^\top A y^*
= v^* = (x^*)^\top A y^*$, so $x^*$ is a best response to $y^*$. Similarly, $y^*$
is a best response to $x^*$. Hence $(x^*, y^*)$ is a NE.

\textbf{(b)} Follows from part (a): every NE strategy pair achieves the value $v^*$.

\textbf{(c)} If $x^*$ and $x^{**}$ are both maximin and $y^*$, $y^{**}$ both minimax,
then by part (a), any combination of a maximin and a minimax strategy forms a NE.
\end{proof}

\begin{rlconnection}[Minimax in RL]
The minimax theorem is the theoretical foundation of:
\begin{itemize}[nosep]
    \item \textbf{Minimax-Q learning} (Littman, 1994): at each state of a two-player
    zero-sum stochastic game, the Bellman operator computes
    $V(s) = \max_x \min_y x^\top Q(s, \cdot, \cdot) y$, which is well-defined
    by the minimax theorem.
    \item \textbf{AlphaGo/AlphaZero}: the MCTS search implicitly approximates
    minimax values in the game tree.
    \item \textbf{GANs}: the generator-discriminator game is a continuous zero-sum
    game $\min_G \max_D V(D,G)$. The minimax theorem (extended to continuous
    strategy spaces) justifies seeking saddle points.
\end{itemize}
\end{rlconnection}

\subsection{Connection to Linear Programming}
\label{subsec:lp}

\begin{proposition}[Solving Zero-Sum Games via LP]
\label{prop:lp_zerosum}
The maximin strategy $x^*$ and the game value $v^*$ can be computed by solving the
LP~\eqref{eq:primal}. The minimax strategy $y^*$ is recovered from the dual
LP~\eqref{eq:dual}. Both LPs have polynomial-time algorithms (e.g., interior point
methods), so computing NE in two-player zero-sum games is in~P.
\end{proposition}

\begin{proof}
This follows directly from the reformulation in the proof of
Theorem~\ref{thm:minimax}. The LP~\eqref{eq:primal} has $m + 1$ variables and
$n + m + 1$ constraints (including non-negativity and the simplex constraint), so
it can be solved efficiently.
\end{proof}

\begin{warning}[Zero-Sum vs.\ General-Sum Complexity]
While NE in two-player zero-sum games can be found in polynomial time via LP,
computing a Nash equilibrium in a general two-player game is PPAD-complete
(Chen, Deng, Teng, 2009; Daskalakis, Goldberg, Papadimitriou, 2009). This
complexity-theoretic barrier is a fundamental challenge for MARL: even
\emph{finding} the solution concept that algorithms like Nash-Q target is
computationally intractable in general.
\end{warning}

%% ============================================================
%% SECTION 7: CORRELATED EQUILIBRIUM
%% ============================================================
\section{Correlated Equilibrium}
\label{sec:correlated}

\begin{definition}[Correlated Equilibrium]
\label{def:ce}
A \emph{correlated equilibrium} (CE) of a finite game $\Gamma = (N, (S_i), (u_i))$
is a probability distribution $\pi$ over $S = \prod_{i \in N} S_i$ such that for
every player~$i$ and every pair of strategies $s_i, s_i' \in S_i$:
\[
    \sum_{s_{-i} \in S_{-i}} \pi(s_i, s_{-i})\, u_i(s_i, s_{-i})
    \geq \sum_{s_{-i} \in S_{-i}} \pi(s_i, s_{-i})\, u_i(s_i', s_{-i}).
\]
Equivalently, if a mediator draws $s = (s_1, \ldots, s_n)$ from $\pi$ and
privately recommends $s_i$ to player~$i$, no player has an incentive to deviate
from the recommendation (assuming the others follow theirs).
\end{definition}

\begin{theorem}[Every Nash Equilibrium Is a Correlated Equilibrium]
\label{thm:ne_is_ce}
If $\sigma^* = (\sigma_1^*, \ldots, \sigma_n^*)$ is a mixed strategy NE, then the
product distribution $\pi(s) = \prod_{i \in N} \sigma_i^*(s_i)$ is a correlated
equilibrium.
\end{theorem}

\begin{proof}
Fix player~$i$ and strategies $s_i, s_i' \in S_i$. The CE incentive constraint
requires:
\[
    \sum_{s_{-i}} \pi(s_i, s_{-i})\, u_i(s_i, s_{-i})
    \geq \sum_{s_{-i}} \pi(s_i, s_{-i})\, u_i(s_i', s_{-i}).
\]
Since $\pi$ is a product distribution, $\pi(s_i, s_{-i}) = \sigma_i^*(s_i) \cdot
\prod_{j \neq i} \sigma_j^*(s_j)$. If $\sigma_i^*(s_i) = 0$, both sides are~0
and the inequality holds. If $\sigma_i^*(s_i) > 0$, dividing both sides by
$\sigma_i^*(s_i)$ gives:
\[
    \sum_{s_{-i}} \bigg(\prod_{j \neq i} \sigma_j^*(s_j)\bigg) u_i(s_i, s_{-i})
    \geq \sum_{s_{-i}} \bigg(\prod_{j \neq i} \sigma_j^*(s_j)\bigg) u_i(s_i', s_{-i}),
\]
i.e., $u_i(s_i, \sigma_{-i}^*) \geq u_i(s_i', \sigma_{-i}^*)$.
Since $s_i \in \supp(\sigma_i^*)$, the indifference principle
(Theorem~\ref{thm:indifference}) gives
$u_i(s_i, \sigma_{-i}^*) \geq u_i(s_i', \sigma_{-i}^*)$ for all $s_i'$, which is
exactly what we need.
\end{proof}

\begin{proposition}[CE as a Linear Program]
\label{prop:ce_lp}
The set of correlated equilibria of a finite game is a convex polytope, defined by the
linear incentive-compatibility constraints in Definition~\ref{def:ce} together with
$\pi(s) \geq 0$ for all $s$ and $\sum_s \pi(s) = 1$. In particular, optimizing any
linear objective over the CE set is a linear program.
\end{proposition}

\begin{proof}
The CE constraints are:
\[
    \sum_{s_{-i}} \pi(s_i, s_{-i})\big[u_i(s_i, s_{-i}) - u_i(s_i', s_{-i})\big] \geq 0
    \quad \text{for all } i, s_i, s_i'.
\]
These are linear in the variables $\{\pi(s)\}_{s \in S}$. Combined with the simplex
constraints $\pi(s) \geq 0$ and $\sum_s \pi(s) = 1$, the feasible set is a polytope
(intersection of finitely many half-spaces). Any linear function of $\pi$ can be
optimized over this polytope by solving an LP.
\end{proof}

\begin{example}[Correlated Equilibrium in Battle of the Sexes]
\label{ex:ce_bos}
In the Battle of the Sexes (Example~\ref{ex:bos}), consider the ``traffic light''
correlated equilibrium: a mediator recommends $(O,O)$ with probability $1/2$ and
$(F,F)$ with probability $1/2$. Under this distribution $\pi$:
\[
    \pi(O,O) = 1/2, \quad \pi(O,F) = 0, \quad \pi(F,O) = 0, \quad \pi(F,F) = 1/2.
\]
Check the CE condition for player~1, given recommendation $O$ (so $s_1 = O$):
\[
    \pi(O,O)\, u_1(O,O) = \tfrac{1}{2} \cdot 3 = \tfrac{3}{2}
    \geq \pi(O,O)\, u_1(F,O) = \tfrac{1}{2} \cdot 0 = 0. \quad \checkmark
\]
Given recommendation $F$:
\[
    \pi(F,F)\, u_1(F,F) = \tfrac{1}{2} \cdot 1 = \tfrac{1}{2}
    \geq \pi(F,F)\, u_1(O,F) = \tfrac{1}{2} \cdot 0 = 0. \quad \checkmark
\]
Similarly for player~2. The expected payoffs are
$(1/2)(3) + (1/2)(1) = 2$ for player~1 and $(1/2)(1) + (1/2)(3) = 2$ for player~2,
which Pareto-dominates the mixed NE payoffs of $(3/4, 3/4)$.
\end{example}

\begin{rlconnection}[Correlated Equilibrium and MARL]
Correlated equilibrium has several advantages over Nash equilibrium for MARL:
\begin{itemize}[nosep]
    \item \textbf{Computational tractability}: finding a CE is an LP, polynomial in
    the game size. Finding a NE is PPAD-complete.
    \item \textbf{Natural emergence}: simple no-regret learning dynamics (e.g., regret
    matching) converge to the set of CE, not NE. This makes CE a more natural
    prediction for what independent learners will achieve.
    \item \textbf{Better welfare}: CE can achieve higher social welfare than any NE,
    as the Battle of the Sexes example illustrates.
\end{itemize}
\end{rlconnection}

%% ============================================================
%% SECTION 8: EXTENSIVE-FORM GAMES
%% ============================================================
\section{Extensive-Form Games}
\label{sec:extensive}

\subsection{Game Trees}
\label{subsec:game_trees}

\begin{definition}[Extensive-Form Game]
\label{def:extensive}
An \emph{extensive-form game} is a tuple
$\Gamma = (N, H, P, (\mathcal{I}_i)_{i \in N}, u)$ where:
\begin{enumerate}[nosep,label=(\roman*)]
    \item $N = \{1, \ldots, n\}$ is the set of players. We may include a ``chance''
    player~0 (nature).
    \item $H$ is a set of \emph{histories} (finite sequences of actions), including the
    empty sequence $\emptyset$ (the root). $H$ is closed under prefixes: if
    $(a_1, \ldots, a_k) \in H$ then $(a_1, \ldots, a_l) \in H$ for all $l \leq k$.
    A history $h \in H$ is \emph{terminal} if there is no $a$ with $(h, a) \in H$.
    The set of terminal histories is $Z \subseteq H$.
    \item $P : H \setminus Z \to N \cup \{0\}$ is the \emph{player function}, assigning
    to each non-terminal history the player who moves at that point.
    \item For each player $i \in N$, $\mathcal{I}_i$ is a partition of
    $\{h \in H \setminus Z : P(h) = i\}$ into \emph{information sets}. If
    $h, h'$ are in the same information set $I \in \mathcal{I}_i$, then the
    set of available actions $A(h) = A(h')$ is the same (denoted $A(I)$).
    \item $u = (u_1, \ldots, u_n) : Z \to \R^n$ assigns payoffs to terminal histories.
\end{enumerate}
\end{definition}

\begin{definition}[Perfect Information]
\label{def:perfect_info}
An extensive-form game has \emph{perfect information} if every information set is a
singleton: $|I| = 1$ for all $I \in \mathcal{I}_i$, for all $i \in N$. In a game
with perfect information, every player knows exactly which node the game has reached
when it is their turn to move.
\end{definition}

\begin{definition}[Perfect Recall]
\label{def:perfect_recall}
Player~$i$ has \emph{perfect recall} if, for any two histories $h, h'$ in the same
information set $I \in \mathcal{I}_i$, player~$i$'s sequence of information sets and
actions leading to $h$ and $h'$ is the same. Informally, no player forgets their own
past actions or information.
\end{definition}

\begin{example}[Entry Game]
\label{ex:entry}
An incumbent (player~2) faces a potential entrant (player~1). The entrant first
decides whether to \emph{Enter} ($E$) or \emph{Stay Out} ($O$). If the entrant enters,
the incumbent decides whether to \emph{Fight} ($F$) or \emph{Accommodate} ($A$):
\begin{itemize}[nosep]
    \item $O$: payoffs $(0, 2)$.
    \item $(E, F)$: payoffs $(-1, -1)$.
    \item $(E, A)$: payoffs $(1, 1)$.
\end{itemize}
This is a game with perfect information (all information sets are singletons).
\end{example}

\subsection{Strategies in Extensive-Form Games}
\label{subsec:strategies_extensive}

\begin{definition}[Pure Strategy in Extensive Form]
\label{def:strategy_extensive}
A \emph{pure strategy} for player~$i$ in an extensive-form game is a function
$s_i : \mathcal{I}_i \to \bigcup_{I \in \mathcal{I}_i} A(I)$ satisfying
$s_i(I) \in A(I)$ for each $I \in \mathcal{I}_i$.
\end{definition}

\begin{remark}[Strategies as Complete Contingent Plans]
\label{rem:contingent}
A strategy specifies an action at \emph{every} information set, including those that
cannot be reached given earlier choices. In the Entry Game (Example~\ref{ex:entry}),
player~2's strategy must specify what to do if the entrant enters, even if that
node is never reached. This redundancy is essential for the normal-form representation
and for defining subgame perfection.
\end{remark}

\begin{definition}[Behavioral Strategy]
\label{def:behavioral}
A \emph{behavioral strategy} for player~$i$ is a function
$b_i : \mathcal{I}_i \to \bigcup_{I \in \mathcal{I}_i} \Delta(A(I))$
satisfying $b_i(I) \in \Delta(A(I))$. That is, $b_i$ assigns a probability
distribution over actions at each information set, with independent randomization
at different information sets.
\end{definition}

\begin{proposition}[Normal-Form Representation]
\label{prop:normal_form_rep}
Every extensive-form game has a corresponding normal-form game: the pure strategy
sets are as in Definition~\ref{def:strategy_extensive}, and the payoff
$u_i(s_1, \ldots, s_n)$ is determined by tracing the unique path through the game
tree when each player follows their strategy (resolving chance moves by expectation).
\end{proposition}

\begin{proof}
Given a strategy profile $(s_1, \ldots, s_n)$, the unique path from the root to a
terminal node is determined inductively: at each non-terminal history $h$, the next
action is $s_{P(h)}(I)$ where $I$ is the information set containing $h$ (or a
probability distribution if $P(h) = 0$, i.e., nature moves). Since the game tree is
finite, this process terminates at some terminal history $z \in Z$. Define
$u_i(s_1, \ldots, s_n)$ as the expected value of $u_i(z)$ where the expectation is
over any chance moves. The resulting tuple $(N, (S_i), (u_i))$ is a well-defined
normal-form game.
\end{proof}

\begin{theorem}[Kuhn's Theorem]
\label{thm:kuhn}
In a finite extensive-form game with perfect recall, every mixed strategy of player~$i$
has a \emph{realization-equivalent} behavioral strategy: a behavioral strategy that
induces the same probability distribution over terminal histories for every strategy
profile of the other players. Conversely, every behavioral strategy has a
realization-equivalent mixed strategy.
\end{theorem}

\begin{proof}
We prove the direction from mixed to behavioral, which is the more useful one.

Let $\sigma_i$ be a mixed strategy (a distribution over pure strategies) for player~$i$
in a game with perfect recall. We construct a behavioral strategy $b_i$ as follows.

For each information set $I \in \mathcal{I}_i$ and action $a \in A(I)$, define
\[
    b_i(a | I) = \frac{\sigma_i\big(\{s_i : s_i(I) = a \text{ and }
    s_i \text{ reaches } I\}\big)}
    {\sigma_i\big(\{s_i : s_i \text{ reaches } I\}\big)},
\]
where ``$s_i$ reaches $I$'' means that the pure strategy $s_i$, together with some
strategy profile of the opponents, leads to a history in $I$. In a game with perfect
recall, whether $s_i$ reaches $I$ depends only on $s_i$'s earlier actions, so this
is well-defined. If the denominator is zero (no pure strategy in the support of
$\sigma_i$ reaches $I$), set $b_i(\cdot | I)$ arbitrarily.

The key property of perfect recall is that it makes the randomization at different
information sets conditionally independent: knowing that player~$i$ reaches information
set $I$ does not reveal what $\sigma_i$ prescribes at other information sets that are
not on the path to $I$. This ensures that the behavioral strategy $b_i$ produces the
same realization probabilities (probability of reaching each terminal node) as $\sigma_i$.

The converse (behavioral to mixed) follows by defining $\sigma_i(s_i) = \prod_I b_i(s_i(I)|I)$,
viewing each pure strategy as specifying an action at each information set.
\end{proof}

\begin{rlconnection}[Behavioral Strategies and RL Policies]
A behavioral strategy is exactly an RL policy: it maps information states (information
sets) to probability distributions over actions. Kuhn's theorem justifies the standard
RL assumption that we can restrict attention to behavioral strategies (policies)
without loss of generality. The perfect recall condition is the Markov property in
disguise: the agent's policy can depend on the current information state without
needing to recall the randomization device used to reach that state.
\end{rlconnection}

%% ============================================================
%% SECTION 9: SUBGAME PERFECT EQUILIBRIUM
%% ============================================================
\section{Subgame Perfect Equilibrium}
\label{sec:spe}

\subsection{Motivation: Non-Credible Threats}
\label{subsec:noncredible}

\begin{example}[Non-Credible Threat in the Entry Game]
\label{ex:noncredible}
In the Entry Game (Example~\ref{ex:entry}), consider the strategy profile where
player~1 plays $O$ (stay out) and player~2 plays $F$ (fight). The payoffs are
$(0, 2)$. Is this a NE?

Player~1's payoff from deviating to $E$ (enter): given player~2's strategy $F$,
the payoff is $-1 < 0$. So player~1 has no incentive to deviate.

Player~2's strategy $F$ is not tested because player~1 stays out, so player~2's
information set is not reached. Player~2 trivially has no incentive to deviate
(any strategy gives the same outcome of $(0,2)$ when player~1 plays $O$).

Therefore $(O, F)$ is a NE. But it relies on a \emph{non-credible threat}:
player~2 threatens to fight, but if the entrant actually entered, fighting
gives player~2 a payoff of $-1$ versus $1$ for accommodating. Rational player~2
would never actually fight. The Nash equilibrium concept fails to exclude such
non-credible threats because it does not require rational behavior at unreached
information sets.
\end{example}

\subsection{Subgames and Subgame Perfection}
\label{subsec:subgame_def}

\begin{definition}[Subgame]
\label{def:subgame}
A \emph{subgame} of an extensive-form game $\Gamma$ beginning at history $h \in H \setminus Z$
is the restriction of $\Gamma$ to the histories that extend $h$, provided:
\begin{enumerate}[nosep,label=(\roman*)]
    \item $h$ is a singleton information set: $\{h\}$ is an information set for
    player $P(h)$.
    \item Every information set that contains a successor of $h$ is entirely contained
    within the subtree rooted at $h$.
\end{enumerate}
The subgame includes all histories $h' \in H$ such that $h$ is a prefix of $h'$,
with the same player function, information sets (restricted to the subtree), and payoffs.
\end{definition}

\begin{remark}
In a game with perfect information, every non-terminal history $h$ begins a subgame
(since all information sets are singletons). In games with imperfect information,
subgames may be rare --- the entire game is always a subgame (beginning at the root).
\end{remark}

\begin{definition}[Subgame Perfect Equilibrium]
\label{def:spe}
A strategy profile $s^*$ is a \emph{subgame perfect equilibrium} (SPE) if its
restriction to every subgame is a Nash equilibrium of that subgame.
\end{definition}

\begin{proposition}[SPE Refines NE]
\label{prop:spe_refines}
Every subgame perfect equilibrium is a Nash equilibrium.
\end{proposition}

\begin{proof}
The entire game is itself a subgame (beginning at the root). If $s^*$ is an SPE,
then it is a NE of the entire game by definition.
\end{proof}

\begin{proposition}[SPE in the Entry Game]
\label{prop:spe_entry}
The Entry Game (Example~\ref{ex:entry}) has a unique SPE: $(E, A)$ (Enter, Accommodate)
with payoffs $(1, 1)$.
\end{proposition}

\begin{proof}
The game has two subgames: the entire game (rooted at the empty history) and the
subgame starting after player~1 chooses $E$ (where player~2 decides between $F$ and $A$).

In the subgame after $E$: player~2 chooses between $F$ (payoff $-1$) and $A$ (payoff $1$).
The unique NE of this subgame is $A$.

Given that player~2 plays $A$ in the subgame, player~1 in the full game chooses between
$O$ (payoff $0$) and $E$ (payoff $1$, since player~2 accommodates). Player~1 chooses $E$.

The unique SPE is $(E, A)$. The non-credible equilibrium $(O, F)$ is eliminated because
$F$ is not a NE in the subgame after entry.
\end{proof}

\subsection{Backward Induction}
\label{subsec:backward_induction}

\begin{definition}[Backward Induction]
\label{def:backward_induction}
For a finite game with perfect information, \emph{backward induction} is the following
procedure:
\begin{enumerate}[nosep]
    \item Identify all ``penultimate'' nodes: non-terminal histories $h$ such that all
    successors $(h, a)$ are terminal. At each such $h$, the player $P(h)$ chooses the
    action maximizing their payoff: $a^*(h) = \argmax_{a \in A(h)} u_{P(h)}(h, a)$.
    \item Replace each penultimate node $h$ with a terminal node having payoff
    $u(h, a^*(h))$.
    \item Repeat until the root is reached.
\end{enumerate}
The resulting strategy profile specifies an action at every node of the original tree.
\end{definition}

\begin{theorem}[Backward Induction Yields SPE]
\label{thm:bi_spe}
In a finite game with perfect information and no ties (i.e., the payoff function is
injective when restricted to any player's choices at any node), backward induction
produces the unique subgame perfect equilibrium.
\end{theorem}

\begin{proof}
We prove by strong induction on the depth $d$ of the game tree (the maximum number of
actions in any history).

\textbf{Base case ($d = 1$):} The game consists of a single player choosing an action
at the root, with all successors being terminal. The backward induction strategy
(choose the payoff-maximizing action) is trivially the unique NE and SPE.

\textbf{Inductive step:} Assume the result holds for all games of depth $< d$.
Consider a game $\Gamma$ of depth $d$ with perfect information and no ties.

At the root, player $P(\emptyset)$ chooses action $a \in A(\emptyset)$. Each action $a$
leads to a subgame $\Gamma_a$ of depth $< d$. By the inductive hypothesis, each
$\Gamma_a$ has a unique SPE $s^*_a$ with a well-defined payoff vector $v(a)$.

In the full game, backward induction assigns to the root the action
$a^* = \argmax_{a \in A(\emptyset)} v_{P(\emptyset)}(a)$, which is unique (no ties).

We show the resulting strategy profile $s^*$ is the unique SPE.

\emph{$s^*$ is an SPE}: In any subgame $\Gamma_a$, the restriction of $s^*$ is $s^*_a$,
which is an SPE of $\Gamma_a$ by the inductive hypothesis. It remains to check that no
player wants to deviate at the root. But the root player $P(\emptyset)$ anticipates the
SPE outcome in each subgame (by the inductive hypothesis, the SPE outcome is unique),
and $a^*$ maximizes over these outcomes. The other players do not move at the root.

\emph{Uniqueness}: Any SPE $\hat{s}$ must restrict to the unique SPE in each subgame
$\Gamma_a$ (by the inductive hypothesis), so the root player faces the same payoffs
$v(a)$ and, with no ties, must choose $a^*$. Hence $\hat{s} = s^*$.
\end{proof}

\subsection{Zermelo's Theorem}
\label{subsec:zermelo}

\begin{keyresult}[Zermelo's Theorem]
\begin{theorem}[Zermelo, 1913]
\label{thm:zermelo}
Every finite extensive-form game with perfect information has a pure strategy subgame
perfect equilibrium. If the game also has no ties (each player has distinct payoffs at
all terminal nodes reachable from any decision node), the SPE is unique.
\end{theorem}
\end{keyresult}

\begin{proof}
\textbf{Existence.} We prove by strong induction on the number of non-terminal
histories $|H \setminus Z|$.

\emph{Base case}: If $|H \setminus Z| = 1$, the game has a single decision node (the
root) followed by terminal nodes. The player at the root has at least one action
maximizing their payoff (the strategy set is finite and non-empty). Any such action
constitutes a pure strategy SPE.

\emph{Inductive step}: Let $\Gamma$ have $|H \setminus Z| = k > 1$. The root player
$P(\emptyset)$ has actions $A(\emptyset) = \{a_1, \ldots, a_l\}$, and each action $a_j$
leads to a subgame $\Gamma_{a_j}$ with strictly fewer non-terminal histories. By the
inductive hypothesis, each $\Gamma_{a_j}$ has a pure strategy SPE $s^*_{a_j}$ with
payoff vector $v(a_j) \in \R^n$.

Define the strategy profile $s^*$ as follows: at the root, choose
$a^* \in \argmax_{a_j} v_{P(\emptyset)}(a_j)$ (this set is non-empty since $l$ is
finite). In each subgame $\Gamma_{a_j}$, follow $s^*_{a_j}$.

We verify $s^*$ is an SPE. For any subgame not containing the root, the restriction
of $s^*$ is $s^*_{a_j}$ for the appropriate $j$, which is an SPE by hypothesis.
For the entire game (which is also a subgame): the root player is choosing
optimally given the SPE continuations (by construction), and no other player moves at the
root, so no player has a profitable deviation. Hence $s^*$ is an SPE.

\textbf{Uniqueness under no ties.} This follows from the backward induction argument
in Theorem~\ref{thm:bi_spe}: at each node, the argmax is unique when payoffs are
distinct, so the SPE is unique.
\end{proof}

\begin{corollary}[Chess Has a Determined Outcome]
\label{cor:chess}
In chess (viewed as a finite game with perfect information --- the 50-move rule and
threefold repetition ensure finiteness), exactly one of the following holds: White has
a winning strategy, Black has a winning strategy, or both players can force a draw.
\end{corollary}

\begin{proof}
Chess is a two-player, finite, perfect-information game with three outcomes: White wins,
Black wins, or draw. By Zermelo's theorem, a pure strategy SPE exists. In this SPE,
the outcome is determined: it is either a White win, a Black win, or a draw. The SPE
strategy for the favored player (or both in the draw case) is the ``winning'' (or
``drawing'') strategy that guarantees this outcome regardless of the opponent's play.
\end{proof}

\begin{rlconnection}[Backward Induction and Dynamic Programming]
Backward induction in game theory is the multi-agent analog of dynamic programming in
single-agent MDPs. In both cases, we work backward from terminal states, computing
optimal decisions at each stage given the optimal continuation. The difference is:
\begin{itemize}[nosep]
    \item \textbf{Single-agent DP}: The Bellman equation $V(s) = \max_a [r(s,a) +
    \gamma V(s')]$ has a unique optimal value function (contraction mapping theorem).
    \item \textbf{Multi-agent backward induction}: At each node, the ``optimal''
    action depends on the other player's SPE strategy, making it a game-theoretic
    rather than optimization problem. With perfect information and no ties, the
    SPE is still unique. With imperfect information or ties, multiplicity arises.
\end{itemize}
AlphaGo and AlphaZero use MCTS (Monte Carlo Tree Search) to approximately perform
backward induction in game trees too large for exact computation.
\end{rlconnection}

\subsection{Behavioral Strategies vs.\ Mixed Strategies in Extensive Form}
\label{subsec:behavioral_vs_mixed}

\begin{remark}[Kuhn's Theorem Revisited]
\label{rem:kuhn_revisited}
Theorem~\ref{thm:kuhn} establishes that in games with perfect recall, we can work
with behavioral strategies (RL policies) without loss of generality. This is
crucial for computational purposes: a behavioral strategy specifies a distribution at
each information set independently, requiring $O(\sum_I |A(I)|)$ parameters, whereas
a mixed strategy over pure strategies may require exponentially many parameters
(one probability for each complete contingent plan).
\end{remark}

\begin{proposition}[Failure Without Perfect Recall]
\label{prop:no_perfect_recall}
In games without perfect recall, mixed and behavioral strategies may yield different
sets of reachable outcome distributions. There exist games without perfect recall
in which a mixed strategy NE exists but no behavioral strategy NE does.
\end{proposition}

\begin{proof}
Consider a one-player game where the player moves twice but forgets their first action
at the second move. Specifically: the player first chooses $L$ or $R$, then (in a
single information set containing both post-$L$ and post-$R$ histories) chooses
$A$ or $B$. Payoffs:
\begin{itemize}[nosep]
    \item $(L, A)$: payoff $1$, \quad $(L, B)$: payoff $0$,
    \item $(R, A)$: payoff $0$, \quad $(R, B)$: payoff $1$.
\end{itemize}
A behavioral strategy specifies $p = \Prob(\text{choose } L)$ at the first node and
$q = \Prob(\text{choose } A)$ at the second information set (which contains both
post-$L$ and post-$R$). The expected payoff is $pq + (1-p)(1-q) = 1 - p - q + 2pq$.
The maximum over behavioral strategies is $1/2$ (achieved, e.g., at $p = q = 1/2$).

However, the mixed strategy that plays $(L, A)$ with probability $1/2$ and $(R, B)$
with probability $1/2$ achieves expected payoff $1/2 \cdot 1 + 1/2 \cdot 1 = 1 > 1/2$.
This mixed strategy has no realization-equivalent behavioral strategy, since the
behavioral strategy cannot correlate the first and second moves.
\end{proof}

%% ============================================================
%% SECTION 10: EXERCISES
%% ============================================================
\section{Exercises}
\label{sec:exercises}

\begin{exercise}[Pure NE Enumeration]
\label{ex:pure_ne}
Find all pure strategy Nash equilibria of the following game:
\[
\begin{array}{c|ccc}
 & L & C & R \\ \hline
T & (2,1) & (0,4) & (1,2) \\
M & (3,2) & (1,1) & (0,3) \\
B & (1,0) & (2,3) & (3,1)
\end{array}
\]
For each NE, verify the best response condition.
\end{exercise}

\begin{exercise}[IESDS Application]
\label{ex:iesds}
Apply iterated elimination of strictly dominated strategies to the game:
\[
\begin{array}{c|ccc}
 & L & C & R \\ \hline
T & (3,1) & (0,0) & (1,3) \\
M & (1,2) & (2,4) & (3,0) \\
B & (0,3) & (4,2) & (2,1)
\end{array}
\]
At each step, identify the dominated strategy and the strategy that dominates it.
State the surviving strategy profiles.
\end{exercise}

\begin{exercise}[Mixed NE in $2 \times 2$ Games]
\label{ex:mixed_2x2}
Compute all Nash equilibria (pure and mixed) of the Stag Hunt
(Example~\ref{ex:stag}). Verify the indifference principle for the mixed NE.
\end{exercise}

\begin{exercise}[Minimax Computation]
\label{ex:minimax}
For the zero-sum game with payoff matrix
\[
A = \begin{pmatrix} 3 & -1 \\ -2 & 4 \end{pmatrix},
\]
compute the value of the game and the optimal mixed strategies for both players.
Formulate the problem as a linear program and verify the minimax theorem directly.
\end{exercise}

\begin{exercise}[Correlated Equilibrium Construction]
\label{ex:ce_construct}
Consider the game of Chicken:
\[
\begin{array}{c|cc}
 & D & S \\ \hline
D & (-5,-5) & (2,-1) \\
S & (-1,2) & (0,0)
\end{array}
\]
\begin{enumerate}[nosep,label=(\alph*)]
    \item Find all pure and mixed Nash equilibria.
    \item Construct a correlated equilibrium that Pareto-dominates the mixed NE.
    \item Verify that your CE satisfies the incentive-compatibility constraints.
\end{enumerate}
\end{exercise}

\begin{exercise}[Backward Induction]
\label{ex:backward}
Consider the following three-stage perfect-information game: Player~1 chooses $A$ or $B$.
If $A$, player~2 chooses $C$ or $D$. If $B$, player~2 chooses $E$ or $F$.
If $D$, player~1 chooses $G$ or $H$. Terminal payoffs (player~1, player~2):
\begin{itemize}[nosep]
    \item $A \to C$: $(3,2)$, \quad $A \to D \to G$: $(4,1)$, \quad $A \to D \to H$: $(1,3)$,
    \item $B \to E$: $(2,4)$, \quad $B \to F$: $(5,0)$.
\end{itemize}
Apply backward induction to find the SPE. List the complete strategy for each player.
\end{exercise}

\begin{exercise}[Non-Credible Threats]
\label{ex:noncredible}
Consider a modified entry game: an entrant chooses $E$ or $O$. If $E$, the incumbent
chooses $F$, $A$, or $P$ (where $P$ means ``price war''). Payoffs:
\begin{itemize}[nosep]
    \item $O$: $(0, 3)$, \quad $(E, F)$: $(-2, -1)$, \quad $(E, A)$: $(1, 1)$, \quad $(E, P)$: $(-3, -2)$.
\end{itemize}
\begin{enumerate}[nosep,label=(\alph*)]
    \item Find all pure strategy NE of the normal-form representation.
    \item Identify which NE involve non-credible threats.
    \item Find the unique SPE.
\end{enumerate}
\end{exercise}

\begin{exercise}[Kuhn's Theorem Application]
\label{ex:kuhn}
In a two-player signaling game, player~1 (the sender) observes a private type
$t \in \{H, L\}$ (each equally likely) and sends a message $m \in \{h, l\}$.
Player~2 (the receiver) observes the message but not the type, and chooses an action
$a \in \{U, D\}$. Player~1 has four pure strategies (mapping types to messages).
\begin{enumerate}[nosep,label=(\alph*)]
    \item How many pure strategies does player~1 have? Player~2?
    \item Write down the normal-form representation.
    \item Verify that this game has perfect recall.
    \item Construct a behavioral strategy for player~1 and its realization-equivalent
    mixed strategy.
\end{enumerate}
\end{exercise}

\begin{exercise}[Properties of NE]
\label{ex:ne_properties}
Prove or disprove the following statements for finite games:
\begin{enumerate}[nosep,label=(\alph*)]
    \item Every game with a dominant strategy equilibrium has a unique NE.
    \item If IESDS yields a unique strategy profile, it is a NE.
    \item The set of NE is always convex.
    \item In a two-player zero-sum game, if $(x^*, y^*)$ and $(x^{**}, y^{**})$ are
    both NE, then $(x^*, y^{**})$ is also a NE.
\end{enumerate}
\end{exercise}

\begin{exercise}[SPE Existence Without Perfect Information]
\label{ex:spe_imperfect}
Construct a two-player extensive-form game with imperfect information that has:
\begin{enumerate}[nosep,label=(\alph*)]
    \item Exactly one subgame (the entire game).
    \item A NE that is not ``sequentially rational'' at the non-singleton information set.
\end{enumerate}
This illustrates why SPE has limited refinement power in games with imperfect information,
motivating stronger concepts like sequential equilibrium.
\end{exercise}

%% ============================================================
%% SECTION 11: SUMMARY
%% ============================================================
\section{Summary of Key Results}
\label{sec:summary}

\begin{center}
\renewcommand{\arraystretch}{1.4}
\begin{tabular}{p{4.5cm} p{8.5cm} l}
\toprule
\textbf{Result} & \textbf{Statement} & \textbf{Ref.} \\
\midrule
Order Independence of IESDS &
In finite games, the set of strategies surviving iterated elimination of strictly
dominated strategies is independent of the elimination order. &
Thm.~\ref{thm:iesds_order} \\

NE as Mutual Best Response &
$s^*$ is a NE iff $s_i^* \in \BR_i(s_{-i}^*)$ for all~$i$. &
Prop.~\ref{prop:ne_br} \\

Indifference Principle &
In a mixed NE, all pure strategies in the support yield equal expected payoff;
strategies outside the support yield weakly less. &
Thm.~\ref{thm:indifference} \\

Nash's Existence Theorem &
Every finite game has at least one mixed strategy NE.
Proved via Kakutani's fixed point theorem. &
Thm.~\ref{thm:nash_existence} \\

Kakutani's Fixed Point Thm &
A uhc correspondence with non-empty convex values on a compact convex set
has a fixed point. Proved via Brouwer. &
Thm.~\ref{thm:kakutani} \\

Von Neumann's Minimax Thm &
$\max_x \min_y x^\top A y = \min_y \max_x x^\top A y$ for any matrix~$A$.
Proved via LP duality. &
Thm.~\ref{thm:minimax} \\

NE in Zero-Sum Games &
NE strategies are interchangeable and all yield the game value. &
Cor.~\ref{cor:ne_zero_sum} \\

NE $\subseteq$ CE &
Every NE (as a product distribution) is a correlated equilibrium. &
Thm.~\ref{thm:ne_is_ce} \\

CE via LP &
The CE set is a convex polytope; optimizing over it is an LP. &
Prop.~\ref{prop:ce_lp} \\

Kuhn's Theorem &
With perfect recall, mixed and behavioral strategies are realization-equivalent. &
Thm.~\ref{thm:kuhn} \\

SPE Refines NE &
Every subgame perfect equilibrium is a Nash equilibrium. &
Prop.~\ref{prop:spe_refines} \\

Backward Induction $\Rightarrow$ SPE &
In finite perfect-information games with no ties, backward induction yields
the unique SPE. &
Thm.~\ref{thm:bi_spe} \\

Zermelo's Theorem &
Every finite perfect-information game has a pure strategy SPE.
Unique if no ties. &
Thm.~\ref{thm:zermelo} \\
\bottomrule
\end{tabular}
\end{center}

\bigskip

\noindent\textbf{Looking ahead.} This lesson establishes the static and single-shot
extensive-form game theory needed for multi-agent RL. The next game theory topics ---
repeated games, stochastic games, and mechanism design --- build directly on these
foundations. In Phase~02 (Core RL), the MDP framework can be viewed as a single-agent
``game against nature.'' In Phase~05 (MARL), stochastic games combine MDPs with the
multi-player game-theoretic structure developed here: at each state, agents play a
stage game whose solution concept (Nash, minimax, correlated) determines the Bellman
operator for the multi-agent value function.

\end{document}
